\documentclass{relprog}

% ─── Metadados do relatório ────────────────────────────────────────────────────
\consultoria{Nome da Consultoria}
\logoConsultoria{logo.png}
\documentoId{REL-PROG-001}
\tituloRelatorio{Relatório de Progresso}
\projetoNome{[Nome do Projeto]}
\clienteNome{[Nome do Cliente]}
\responsavel{[Engenheiro Responsável]}
\revisao{Rev A}
\periodo{[dd/mm/aaaa] a [dd/mm/aaaa]}
\faseAtual{[Nome da Fase]}

% ─── Assinaturas ───────────────────────────────────────────────────────────────
\assinaturaResponsavel{[Nome do Responsável]}{[Cargo]}
\assinaturaCliente{[Nome do Representante do Cliente]}{[Cargo] \textbar\ [Empresa do Cliente]}

\begin{document}

\fazerCapa


%% ════════════════════════════════════════════════════════════════
%%  RESUMO EXECUTIVO
%% ════════════════════════════════════════════════════════════════
\begin{resumobox}
O projeto encontra-se \textbf{[no prazo / com atraso de X dias / adiantado]}.
As principais atividades do período foram \textbf{[descrever em 1 frase]}.
O risco mais relevante no momento é \textbf{[descrever ou escrever ``nenhum identificado'']}.
Há \textbf{[N decisões / nenhuma decisão]} pendente por parte do cliente.
\end{resumobox}


%% ════════════════════════════════════════════════════════════════
%%  1. STATUS POR ÁREA
%% ════════════════════════════════════════════════════════════════
\section{1. Status por Área}

\begin{tabularx}{\linewidth}{@{} l c X @{}}
  \cabStatusArea
  Hardware — Esquemático  & \verde    & Revisão 2 concluída. Aguarda aprovação para layout. \\[3pt]
  Hardware — PCB Layout   & \amarelo  & Em andamento. Roteamento 60\% concluído. \\[3pt]
  Firmware — Drivers BSP  & \verde    & Drivers de SPI, I2C e UART validados em bancada. \\[3pt]
  Firmware — Aplicação    & \amarelo  & Módulo de comunicação em desenvolvimento. \\[3pt]
  Testes \& Validação     & \vermelho & Aguarda protótipo físico para iniciar. \\[3pt]
  Compras / Componentes   & \verde    & BOM aprovada. Pedido realizado em [data]. \\
  \bottomrule
\end{tabularx}

\legendaSemaforo


%% ════════════════════════════════════════════════════════════════
%%  2. ATIVIDADES DO PERÍODO
%% ════════════════════════════════════════════════════════════════
\section{2. Atividades do Período}

\subsection{Concluídas}

\begin{listaItens}
  \item \textbf{[Área]:} Descrição objetiva da atividade concluída.
  \item \textbf{[Área]:} Descrição objetiva da atividade concluída.
  \item \textbf{[Área]:} Descrição objetiva da atividade concluída.
\end{listaItens}

\subsection{Em Andamento}

\begin{listaItens}
  \item \textbf{[Área]:} Descrição da atividade em andamento. Previsão: \textbf{[data]}.
  \item \textbf{[Área]:} Descrição da atividade em andamento. Previsão: \textbf{[data]}.
\end{listaItens}

\subsection{Não Iniciadas / Atrasadas}

\begin{listaItens}
  \item \textbf{[Área]:} Descrição. Motivo do atraso: \textbf{[justificativa]}. Nova previsão: \textbf{[data]}.
\end{listaItens}


%% ════════════════════════════════════════════════════════════════
%%  3. PROGRESSO DO CRONOGRAMA
%% ════════════════════════════════════════════════════════════════
\section{3. Progresso do Cronograma}

\begin{tabularx}{\linewidth}{@{} l l l c X @{}}
  \cabCronograma
  Aprovação de Requisitos  & 05/05/2025 & 05/05/2025 & 100\% & Concluído conforme planejado. \\[3pt]
  Esquemático Rev A        & 20/05/2025 & 22/05/2025 &  90\% & 2 dias de atraso por revisão de requisitos. \\[3pt]
  PCB Layout               & 10/06/2025 & 10/06/2025 &  60\% & Em andamento. \\[3pt]
  Firmware v0.1            & 15/06/2025 & 15/06/2025 &  40\% & Em andamento. \\[3pt]
  Protótipo Físico         & 30/06/2025 & 30/06/2025 &   0\% & Aguarda PCB. \\[3pt]
  Testes \& Validação      & 20/07/2025 & 20/07/2025 &   0\% & Aguarda protótipo. \\[3pt]
  Entrega Final            & 31/07/2025 & 31/07/2025 &   0\% & — \\
  \bottomrule
\end{tabularx}


%% ════════════════════════════════════════════════════════════════
%%  4. RISCOS E IMPEDIMENTOS
%% ════════════════════════════════════════════════════════════════
\section{4. Riscos e Impedimentos}

\begin{tabularx}{\linewidth}{@{} l l l X X @{}}
  \cabRiscos
  R1 & Lead time de componentes & Alto    & Médio   & Pedido antecipado realizado. \\[3pt]
  R2 & Mudança de requisitos    & Médio   & Baixo   & Change log formal em vigor. \\[3pt]
  R3 & [Descrever risco]        & [Nível] & [Nível] & [Ação de mitigação]. \\
  \bottomrule
\end{tabularx}


%% ════════════════════════════════════════════════════════════════
%%  5. PRÓXIMOS PASSOS
%% ════════════════════════════════════════════════════════════════
\section{5. Próximos Passos}

\subsection{Ações da Consultoria — próximo período}

\begin{listaItens}
  \item \textbf{[Área]:} Descrição da próxima atividade planejada. Prazo: \textbf{[data]}.
  \item \textbf{[Área]:} Descrição da próxima atividade planejada. Prazo: \textbf{[data]}.
  \item \textbf{[Área]:} Descrição da próxima atividade planejada. Prazo: \textbf{[data]}.
\end{listaItens}

\subsection{Pendências com o Cliente}

\begin{tabularx}{\linewidth}{@{} c X l l @{}}
  \cabPendencias
  P1 & Aprovação do esquemático Rev B para liberar PCB.    & \today & [data] \\[3pt]
  P2 & Confirmação de requisito de consumo em sleep mode.  & \today & [data] \\[3pt]
  P3 & [Descrever pendência do cliente]                    & \today & [data] \\
  \bottomrule
\end{tabularx}

\alertaPendencias


%% ════════════════════════════════════════════════════════════════
%%  6. ASSINATURAS
%% ════════════════════════════════════════════════════════════════
\section{6. Emissão e Aprovação}

\fazerAssinaturas

\end{document}
